% Gabarit pandoc des corrigés publiés (brut et guidé). Contenu sans marque ni nom de site.
% Compilation : lualatex (latexmk -lualatex).
% Le Markdown reste la source unique : ne pas éditer ce fichier .tex à la main.
\documentclass[11pt,a4paper]{article}

% --- Moteur, polices (Latin Modern) et langue ---------------------------------
\usepackage{iftex}
\ifLuaTeX\else\errmessage{Ce document se compile avec lualatex}\fi
\usepackage{amsmath,amssymb,mathtools}
\usepackage{fontspec}
% mathrm=sym : \mathrm prend ses glyphes dans la police mathématique (ℓ, μ, Ω), comme KaTeX.
\usepackage[math-style=TeX,bold-style=TeX,mathrm=sym]{unicode-math}
% Repli pour les glyphes absents de Latin Modern (indices et exposants Unicode,
% flèches, filets des schémas en texte, symboles) : DejaVu, puis Latin Modern Math
% (ℓ, lettres mathématiques Unicode des légendes).
\directlua{
  luaotfload.add_fallback("repliTexte", {"DejaVuSerif:mode=harf;", "DejaVuSans:mode=harf;", "file:latinmodern-math.otf:mode=harf;"})
  luaotfload.add_fallback("repliMono", {"DejaVuSansMono:mode=harf;", "DejaVuSans:mode=harf;", "file:latinmodern-math.otf:mode=harf;"})
}
\setmainfont{Latin Modern Roman}[RawFeature={fallback=repliTexte}]
\setsansfont{Latin Modern Sans}[RawFeature={fallback=repliTexte}]
\setmonofont{Latin Modern Mono}[RawFeature={fallback=repliMono}]
\setmathfont{Latin Modern Math}
\usepackage[french]{babel}
\frenchsetup{SmallCapsFigTabCaptions=false}

% --- Mise en page --------------------------------------------------------------
\usepackage[a4paper,margin=2.2cm,headheight=14pt,footskip=1.2cm]{geometry}
\usepackage{microtype}
\setlength{\emergencystretch}{4em}
\tolerance=1500
\setlength{\parindent}{0pt}
\setlength{\parskip}{0.55\baselineskip plus 2pt minus 1pt}
\usepackage{xcolor}
% Gris nommés : babel-french rend « ! » actif dans le document.
\definecolor{gris2}{gray}{0.98}
\definecolor{gris3}{gray}{0.97}
\definecolor{gris5}{gray}{0.95}
\definecolor{gris35}{gray}{0.65}
\definecolor{gris40}{gray}{0.60}
\definecolor{gris45}{gray}{0.55}
\definecolor{gris50}{gray}{0.50}
\definecolor{gris55}{gray}{0.45}
\definecolor{gris60}{gray}{0.40}
\definecolor{gris65}{gray}{0.35}
\definecolor{gris70}{gray}{0.30}
\definecolor{gris80}{gray}{0.20}
\definecolor{gris85}{gray}{0.15}
\usepackage{graphicx}
\usepackage{float}
\usepackage{caption}
\captionsetup{font=small,labelformat=empty,justification=centering}
\usepackage{booktabs,longtable,array,multirow,calc}
\usepackage{siunitx}
\sisetup{locale=FR}
\usepackage[version=4]{mhchem}
\usepackage{listings}
\usepackage{fancyhdr}
\usepackage{lastpage}
\usepackage[most]{tcolorbox}
\usepackage{enumitem}
\setlist{nosep,leftmargin=*,topsep=2pt}
\usepackage[hidelinks,unicode]{hyperref}
\usepackage{bookmark}
\hypersetup{pdftitle={Mines-Ponts MP-MPI-PC-PSI — Français 2026},pdfauthor={Rédigé par IA},pdfcreator={LaTeX},pdfsubject={Corrigé guidé}}
\usepackage{xurl}

% --- Pandoc : listes, tableaux, figures ----------------------------------------
\providecommand{\passthrough}[1]{#1}
\providecommand{\tightlist}{\setlength{\itemsep}{0pt}\setlength{\parskip}{0pt}}

\makeatletter
\def\maxwidth{\ifdim\Gin@nat@width>\linewidth\linewidth\else\Gin@nat@width\fi}
\def\maxheight{\ifdim\Gin@nat@height>0.55\textheight 0.55\textheight\else\Gin@nat@height\fi}
\makeatother
\setkeys{Gin}{width=\maxwidth,height=\maxheight,keepaspectratio}
\setcounter{secnumdepth}{-\maxdimen}
\providecommand{\pandocbounded}[1]{#1}
\usepackage{longtable,booktabs,array}

% --- Code (listings) -----------------------------------------------------------
\lstset{
  basicstyle=\ttfamily\small,
  columns=fullflexible,
  keepspaces=true,
  breaklines=true,
  breakatwhitespace=false,
  postbreak=\mbox{\textcolor{gris50}{$\hookrightarrow$}\space},
  frame=single,
  rulecolor=\color{gris35},
  framesep=4pt,
  xleftmargin=5pt,xrightmargin=5pt,
  backgroundcolor=\color{gris3},
  keywordstyle=\bfseries,
  commentstyle=\color{gris65},
  stringstyle=\color{gris80},
  showstringspaces=false,
  tabsize=4,
  extendedchars=true,
  aboveskip=6pt,belowskip=6pt,
}

% --- Macros des corrigés (KaTeX ↔ LaTeX) ---------------------------------------
% Les corrigés sont écrits pour KaTeX : on définit ici ce que KaTeX connaît et
% que LaTeX n'a pas, sans jamais retoucher le Markdown.
\providecommand{\R}{\mathbb{R}}
\providecommand{\N}{\mathbb{N}}
\providecommand{\Z}{\mathbb{Z}}
\providecommand{\Q}{\mathbb{Q}}
\providecommand{\C}{\mathbb{C}}
\providecommand{\dd}{\mathrm{d}}
\providecommand{\llbracket}{\lBrack}
\providecommand{\rrbracket}{\rBrack}
\providecommand{\lBrace}{\lbrace}
\providecommand{\rBrace}{\rbrace}
\providecommand{\Box}{\square}
\providecommand{\Diamond}{\lozenge}
\providecommand{\varnothing}{\emptyset}
\providecommand{\argmax}{\operatorname*{arg\,max}}
\providecommand{\argmin}{\operatorname*{arg\,min}}
\renewcommand{\leq}{\leqslant}
\renewcommand{\geq}{\geqslant}

% Équation hors texte : réduite à la largeur de la ligne si elle déborde.
\newsavebox{\boiteequation}
\newcommand{\equationajustee}[1]{%
  \sbox{\boiteequation}{$\displaystyle #1$}%
  \ifdim\wd\boiteequation>\linewidth
    \[\resizebox{\linewidth}{!}{\usebox{\boiteequation}}\]%
  \else
    \[#1\]%
  \fi}

% --- Encadrés pédagogiques (sobres, lisibles en noir et blanc) -----------------
\tcbset{
  encadre/.style={
    enhanced,breakable,sharp corners,
    colback=white,colframe=gris55,
    boxrule=0.4pt,leftrule=2.2pt,
    left=6pt,right=6pt,top=3pt,bottom=3pt,
    before skip=6pt,after skip=6pt,
    fonttitle=\bfseries\small,coltitle=black,
    attach title to upper,after title={\quad},
    parbox=false,
  }
}
\newtcolorbox{indication}[1]{encadre,colframe=gris45,colback=gris2,title={Indication~#1}}
\newtcolorbox{methode}{encadre,colframe=gris70,title={Méthode}}
\newtcolorbox{piege}{encadre,colframe=gris85,boxrule=0.8pt,leftrule=2.2pt,title={Piège}}
\newtcolorbox{jury}{encadre,colframe=gris60,colback=gris5,title={Jury}}
\newtcolorbox{encadrenote}{encadre,colframe=gris40}

% --- Titres de section ---------------------------------------------------------
\usepackage{titlesec}
\titleformat{\section}{\Large\bfseries}{}{0pt}{}
\titleformat{\subsection}{\large\bfseries}{}{0pt}{}
\titleformat{\subsubsection}{\normalsize\bfseries}{}{0pt}{}
\titlespacing*{\section}{0pt}{14pt plus 4pt minus 2pt}{6pt}
\titlespacing*{\subsection}{0pt}{10pt plus 3pt minus 2pt}{4pt}
\titlespacing*{\subsubsection}{0pt}{8pt plus 2pt minus 2pt}{2pt}

% --- En-tête et pied de page ---------------------------------------------------
\pagestyle{fancy}
\fancyhf{}
\fancyhead[L]{\small\textsc{Corrigé guidé} · Mines-Ponts MP-MPI-PC-PSI}
\fancyhead[R]{\small Français 2026}
\fancyfoot[L]{\footnotesize Corrigé rédigé par IA, non officiel}
\fancyfoot[R]{\footnotesize \thepage\,/\,\pageref*{LastPage}}
\renewcommand{\headrulewidth}{0.4pt}
\renewcommand{\footrulewidth}{0pt}
\fancypagestyle{plain}{\fancyhf{}\fancyfoot[R]{\footnotesize \thepage\,/\,\pageref*{LastPage}}\fancyfoot[L]{\footnotesize Corrigé rédigé par IA, non officiel}\renewcommand{\headrulewidth}{0pt}}

\begin{document}
\thispagestyle{plain}
\begin{center}
  {\small\scshape Corrigé guidé}\par\smallskip
  {\LARGE\bfseries Mines-Ponts MP-MPI-PC-PSI — Français 2026\par}\medskip
  {\large Mines-Ponts · MP-MPI-PC-PSI · Français · 2026\par}
\end{center}
\vspace{2pt}\hrule\vspace{8pt}

\emph{Corrigé rédigé par IA, non officiel. Signalez toute erreur.}

\hypertarget{pruxe9sentation-du-sujet}{%
\section{Présentation du sujet}\label{pruxe9sentation-du-sujet}}

\textbf{L'épreuve.} Trois heures, aucun document ni appareil autorisé, épreuve commune à toutes les filières. Un seul exercice : une dissertation sur une citation, à traiter à partir des trois œuvres du thème « Expériences de la nature ». La page de garde demande d'écrire une ligne sur deux. Aucune longueur n'est imposée ; une très bonne copie tient en général entre 2 000 et 3 000 mots, et c'est la qualité de l'argumentation, non le volume, qui fait la note.

\textbf{La citation.} George Sand, « La Forêt de Fontainebleau », article du \emph{Temps} du 13 novembre 1872. Le contexte compte : Sand y défend une forêt exploitée, dont une partie avait été mise en réserve à la demande des peintres. La phrase est une exclamation en quatre temps : une annonce (« Quelle conquête à entreprendre »), trois infinitifs qui forment un programme (entrer, chercher, rapporter) et une longue relative qui en donne le bénéfice (« la possession de son être »). Trois lexiques s'y superposent :

\begin{longtable}[]{@{}
  >{\raggedright\arraybackslash}p{(\columnwidth - 4\tabcolsep) * \real{0.1730}}
  >{\raggedright\arraybackslash}p{(\columnwidth - 4\tabcolsep) * \real{0.3820}}
  >{\raggedright\arraybackslash}p{(\columnwidth - 4\tabcolsep) * \real{0.4149}}@{}}
\toprule\noalign{}
\begin{minipage}[b]{\linewidth}\raggedright
Lexique
\end{minipage} & \begin{minipage}[b]{\linewidth}\raggedright
Mots
\end{minipage} & \begin{minipage}[b]{\linewidth}\raggedright
Ce qu'il suppose
\end{minipage} \\
\midrule\noalign{}
\endhead
\bottomrule\noalign{}
\endlastfoot
la prise & conquête, entreprendre, possession & la nature comme territoire à gagner \\
le sacré & oracle, forêt sacrée, charme & une nature qui parle et demande interprétation \\
l'être & tout homme vivant ou à naître, son être & un enjeu universel, transmis, et intérieur \\
\end{longtable}

Le point décisif, qu'il faut voir au brouillon : \textbf{le complément de la conquête n'est pas la nature}. On entre dans la nature, mais ce qu'on possède à la fin, c'est « son être ». La restrictive « ne fût-ce qu'un mot » tempère par ailleurs l'enthousiasme : le butin est minuscule.

\textbf{Les œuvres.}

\begin{itemize}
\tightlist
\item
  \emph{Vingt mille lieues sous les mers} (Verne, 1869-1870) : le professeur Aronnax, son domestique Conseil et le harponneur Ned Land, prisonniers du capitaine Nemo à bord du \emph{Nautilus}. Exploration, inventaire du vivant marin, mais aussi chasses, conquête du pôle Sud, vengeance de Nemo et fin ouverte dans le Maelström.
\item
  \emph{La Connaissance de la vie} (Canguilhem, 1952 ; édition augmentée de 1965), dans les limites fixées par l'énoncé : l'introduction « La pensée et le vivant », « L'expérimentation en biologie animale » (partie « Méthode ») et les chapitres II à V de la partie « Philosophie ». Thèses utiles : la pensée du vivant procède du vivant ; le vivant est une totalité qu'on n'isole pas sans l'altérer ; la machine s'explique par l'organisme plutôt que l'inverse ; chaque vivant a son milieu propre ; la santé est normativité.
\item
  \emph{Le Mur invisible} (Haushofer, 1963) : une femme, jamais nommée, enfermée par une paroi transparente dans une vallée des Alpes autrichiennes avec un chien (Lynx), une vache (Bella), puis des chats ; elle écrit après coup son « rapport ».
\end{itemize}

\textbf{Difficulté.} Le corpus est hétérogène (roman d'aventures scientifiques, épistémologie, roman de survie) et la citation, lyrique et romantique, semble épouser Verne et rebuter Canguilhem. Le rapport 2026 juge le sujet « très littéraire » : il invitait à réfléchir au langage, à la force du mot et à l'acte de nommer, et il a « déstabilisé bien des candidats ». Le rapport ne donne aucune statistique, mais il désigne clairement ce qui a départagé les copies :

\begin{itemize}
\tightlist
\item
  l'\textbf{analyse de toute la citation}, négligée cette année plus que d'habitude : beaucoup de copies n'ont retenu que « conquête », lue hors contexte, et ont oublié l'oracle, le mot et la possession de son être ;
\item
  une \textbf{problématique unique} qui fasse voir une tension, au lieu d'une liste de questions ;
\item
  une \textbf{troisième partie} qui dépasse vraiment les deux premières, point le plus souvent manqué, même dans les bonnes copies.
\end{itemize}

S'y ajoutent deux exigences propres au corpus : faire de Canguilhem un véritable interlocuteur et non une caution, et exploiter le motif du « rapport », littéralement présent dans les trois œuvres.

Le rapport contient en annexe un exemple complet de dissertation, sur un autre plan que celui proposé ici (conquête de soi ou conquête de la nature, puis cohabiter plutôt que conquérir). Il précise lui-même que ce corrigé excède ce qu'on attend un jour de concours.

\textbf{Gestion du temps conseillée.} Environ 45 minutes pour l'analyse, la problématique et le plan détaillé ; l'introduction rédigée au brouillon (15 minutes) ; environ 1 h 45 pour les trois parties ; 10 minutes pour la conclusion ; un dernier quart d'heure de relecture (orthographe, accords, noms propres, exactitude de la citation).

\hypertarget{premiuxe8re-phase-pruxe9parer-au-brouillon}{%
\section{Première phase --- Préparer au brouillon}\label{premiuxe8re-phase-pruxe9parer-au-brouillon}}

Trois étapes qui ne se voient pas sur la copie mais qui décident de tout le reste : lire la citation mot à mot, en tirer une problématique, bâtir un plan où chaque œuvre sert dans chaque partie.

\hypertarget{uxe9tape-1-analyser-la-citation}{%
\subsection{Étape 1 --- Analyser la citation}\label{uxe9tape-1-analyser-la-citation}}

\begin{indication}{1}

Commencez par la forme : comment la phrase est-elle construite, et où tombe son poids ?

\end{indication}

\begin{indication}{2}

Interrogez chaque mot porteur sur ce qu'il affirme et sur ce qu'il présuppose. Ne négligez pas les petits mots : « et je dis », « ne fût-ce que », « doit », le possessif « son ».

\end{indication}

\begin{indication}{3}

Demandez-vous quel est l'objet de la « conquête » au début de la phrase, puis quel est l'objet de la « possession » à la fin.

\end{indication}

\hypertarget{solution}{%
\subsubsection{Solution}\label{solution}}

La phrase s'ouvre sur un mot de prise, « conquête », et se ferme sur un mot d'intériorité, « son être ». Entre les deux, trois infinitifs marquent trois moments : \textbf{entrer} (franchir un seuil, s'immerger), \textbf{chercher l'oracle} (interroger une nature qui parle mais dont la parole est obscure), \textbf{rapporter le mot} (revenir et transmettre). Le dernier membre, de loin le plus long, n'est plus une action mais un effet : le mot rapporté doit « répandre sur toute sa vie le charme profond de la possession de son être ». Le poids de la période tombe donc sur ce bénéfice final, et non sur la conquête qui l'annonce.

Ce que les mots présupposent fournit les objections :

\begin{longtable}[]{@{}
  >{\raggedright\arraybackslash}p{(\columnwidth - 4\tabcolsep) * \real{0.2657}}
  >{\raggedright\arraybackslash}p{(\columnwidth - 4\tabcolsep) * \real{0.3159}}
  >{\raggedright\arraybackslash}p{(\columnwidth - 4\tabcolsep) * \real{0.3884}}@{}}
\toprule\noalign{}
\begin{minipage}[b]{\linewidth}\raggedright
Terme
\end{minipage} & \begin{minipage}[b]{\linewidth}\raggedright
Ce qu'il affirme
\end{minipage} & \begin{minipage}[b]{\linewidth}\raggedright
Ce qu'il présuppose
\end{minipage} \\
\midrule\noalign{}
\endhead
\bottomrule\noalign{}
\endlastfoot
conquête, entreprendre & une expérience active, difficile, choisie & qu'entrer, c'est prendre ; qu'on y va de son plein gré \\
tout homme vivant ou à naître & une portée universelle et transmise & que l'expérience soit communicable et qu'on doive la préserver \\
entrer dans la nature & l'immersion, non la contemplation & que l'homme se tienne ordinairement au dehors \\
l'oracle de la forêt sacrée & une nature qui parle, un savoir révélé & qu'il y ait un sens à recueillir \\
rapporter le mot & un retour, une transmission & qu'on revienne et qu'on ait un destinataire \\
ne fût-ce qu'un mot & un butin minime mais suffisant & une modestie que les œuvres peuvent confirmer \\
la possession de son être & le vrai gain est intérieur & que se posséder ait un sens autre que posséder un bien \\
\end{longtable}

\begin{methode}

Sur une citation longue, distinguez ce que la phrase affirme de ce qu'elle présuppose : les affirmations nourrissent la première partie, les présupposés la deuxième, et le terme le plus riche (ici « son être ») la troisième.

\end{methode}

\begin{piege}

Lire « conquête » comme un éloge de la domination technique de la nature. Sand défend une forêt menacée, et l'objet final de la conquête est soi-même : un devoir qui s'indigne pendant trois pages contre la « conquête » réfute une thèse que la citation ne soutient pas. Le jury a vu beaucoup de plans recyclés sur la tension domination de la nature contre domination par la nature : ils passent à côté du sujet.

\end{piege}

\begin{jury}

Rapport 2026 : « Le jury s'est très souvent demandé où avaient bien pu passer « l'oracle de la forêt sacrée », « le mot » à « rapporter » et « le charme profond de la possession de son être ». » Le jury admet qu'on s'interroge sur le mot « conquête », qui évoque un imaginaire guerrier, mais rappelle qu'il n'est pas employé sans contexte. Il conseille une phase d'oubli méthodique des fiches et un travail de synonymie et d'antonymie sur chaque expression obscure, avant de convoquer les œuvres.

\end{jury}

\begin{piege}

George Sand est une femme : le rapport 2026 note que beaucoup de candidats l'ont ignoré. Accordez en conséquence (« elle », « l'autrice », « la romancière »).

\end{piege}

\hypertarget{uxe9tape-2-probluxe9matiser}{%
\subsection{Étape 2 --- Problématiser}\label{uxe9tape-2-probluxe9matiser}}

\begin{indication}{1}

La formule « dans quelle mesure\ldots{} éclaire » demande une mesure : jusqu'où la citation vaut, où elle cesse de valoir, à quelle condition on peut la reprendre.

\end{indication}

\begin{indication}{2}

Cherchez la tension entre le mouvement de la phrase (entrer, chercher, rapporter) et son vocabulaire (conquête, oracle, possession). Les œuvres confirment-elles l'un et l'autre au même degré ?

\end{indication}

\hypertarget{solution-1}{%
\subsubsection{Solution}\label{solution-1}}

Les trois œuvres racontent ou pensent bien une immersion, une interrogation et un écrit : le \textbf{mouvement} de la phrase leur convient. Mais son \textbf{vocabulaire} (conquête, oracle, retour, charme) est démenti par chacune : la conquête tourne à la prédation, la nature ne parle pas, le retour est incertain. Reste le dernier terme, qui renverse la phrase de l'intérieur : ce qu'on possède à la fin, c'est son être, non la nature.

Problématique retenue : \emph{l'expérience de la nature, telle que la donnent à lire les trois œuvres, est-elle cette conquête dont l'homme rapporterait un mot qui le rendrait maître de son être, ou bien ne vérifie-t-elle la phrase de Sand qu'à condition d'en renverser le vocabulaire de la prise ?}

\begin{piege}

Deux problématiques à fuir : « la nature est-elle un objet de conquête ? », qui abandonne la citation dès l'introduction, et « Sand a-t-elle raison ? », qui fait le procès d'une autrice hors programme au lieu de relire les œuvres à sa lumière. Le rapport 2026 cite aussi, comme trop vastes et sans tension, « La conquête de la nature permet-elle une vie meilleure ? » ; il déconseille d'ouvrir la problématique par « En quoi » et de la ramener à l'intitulé du thème.

\end{piege}

\begin{jury}

Rapport 2026 : « Le raisonnement doit découler d'une seule problématique, d'une tension sur laquelle se bâtit le plan. » Le jury reproche à beaucoup de copies d'aligner quatre ou cinq questions ; il ajoute : « On réservera donc la modalité interrogative à la problématique. » Il cite en exemple une très bonne copie qui opposait deux voies : entrer dans la nature comme révélation de soi, ou se conquérir soi-même dans l'adversité qu'elle oppose.

\end{jury}

\hypertarget{uxe9tape-3-construire-le-plan-et-ruxe9partir-les-ruxe9fuxe9rences}{%
\subsection{Étape 3 --- Construire le plan et répartir les références}\label{uxe9tape-3-construire-le-plan-et-ruxe9partir-les-ruxe9fuxe9rences}}

\begin{indication}{1}

Un plan par œuvre est à proscrire : il empêche la confrontation. Indexez plutôt chaque partie sur des termes de la citation.

\end{indication}

\begin{indication}{2}

Avant de rédiger, remplissez une grille partie par partie et œuvre par œuvre : chaque case doit contenir une scène ou une thèse précise, et chaque œuvre doit servir dans chaque partie.

\end{indication}

\hypertarget{solution-2}{%
\subsubsection{Solution}\label{solution-2}}

Plan en trois temps, accord, démenti, déplacement :

\begin{itemize}
\tightlist
\item
  \textbf{I. Le trajet de Sand est celui des œuvres.} A. Entrer : quitter le cabinet. B. Chercher l'oracle : interroger une nature qui ne livre pas son sens d'elle-même. C. Rapporter le mot : les trois œuvres sont des rapports, et le savoir revient sur celui qui l'acquiert.
\item
  \textbf{II. Les œuvres récusent le vocabulaire de Sand.} A. La conquête tourne à la prédation. B. L'oracle se tait. C. Le retour est incertain et le mot ne rend pas la vie.
\item
  \textbf{III. Relire la phrase par son dernier mot.} A. On n'entre pas dans la nature, on n'en est jamais sorti. B. « Ne fût-ce qu'un mot » : les œuvres sont ce mot. C. Posséder son être : normativité, soin, responsabilité.
\end{itemize}

Répartition des références :

\begin{longtable}[]{@{}
  >{\raggedright\arraybackslash}p{(\columnwidth - 6\tabcolsep) * \real{0.1406}}
  >{\raggedright\arraybackslash}p{(\columnwidth - 6\tabcolsep) * \real{0.2982}}
  >{\raggedright\arraybackslash}p{(\columnwidth - 6\tabcolsep) * \real{0.2811}}
  >{\raggedright\arraybackslash}p{(\columnwidth - 6\tabcolsep) * \real{0.2501}}@{}}
\toprule\noalign{}
\begin{minipage}[b]{\linewidth}\raggedright
\end{minipage} & \begin{minipage}[b]{\linewidth}\raggedright
Verne
\end{minipage} & \begin{minipage}[b]{\linewidth}\raggedright
Canguilhem
\end{minipage} & \begin{minipage}[b]{\linewidth}\raggedright
Haushofer
\end{minipage} \\
\midrule\noalign{}
\endhead
\bottomrule\noalign{}
\endlastfoot
I.A Entrer & Aronnax, savant de cabinet ; la forêt sous-marine de l'île Crespo & l'expérimentation : le vivant est une totalité, le laboratoire un milieu & l'entrée subie : le mur découvert au matin \\
I.B Chercher & le « monstre » de 1866 ; Conseil classe, Ned reconnaît & le vivant, centre qui organise son milieu & lire le ciel, les bêtes, la neige \\
I.C Rapporter & le récit d'Aronnax ; Aronnax partagé & la pensée du vivant tient du vivant l'idée du vivant & le « rapport » ; la femme qu'elle était devenue étrangère \\
II.A Conquête & le pavillon du pôle Sud puis la banquise ; cachalots ; l'éperonnage & « Machine et organisme » contre le modèle cartésien & l'homme à la hache ; la chasse à contrecœur \\
II.B Oracle & le mot ATLANTIDE écrit par un homme ; l'abîme insondable & « Le vivant et son milieu » : la tique de von Uexküll & le mur inexpliqué ; la corneille blanche \\
II.C Rapporter & Aronnax doute qu'on le croie ; le manuscrit de Nemo confié à la mer ; fin ouverte & connaître, c'est analyser & le papier manque ; pas de destinataire \\
III.A Pas de seuil & scaphandre, \emph{Nautilus}, \emph{Mobilis in mobili} & l'univers sans centre du savant, le milieu centré du vivant & créature parmi d'autres \\
III.B Un mot & « deux hommes » ont le droit de répondre & une manière de penser, non des résultats & écrire pour demeurer quelqu'un \\
III.C Son être & Nemo possède la mer, non lui-même ; Aronnax a changé & la santé comme normativité & normes nouvelles et soin des bêtes \\
\end{longtable}

\begin{jury}

Rapport 2026, sur l'annonce du plan : « Voilà pourquoi nous conseillons vivement de faire une phrase par partie. » Annoncer qu'on va « aller dans le sens de la citation » puis la « nuancer » ne suffit pas : « il faut expliciter l'axe de légitimation de la thèse, tout comme son axe de nuance. Aucun suspense ne doit être ménagé ! »

\end{jury}

\begin{piege}

Sur ce sujet, l'œuvre qu'on oublie est \emph{La Connaissance de la vie} : recueil d'épistémologie sans personnages, elle tend à n'apparaître qu'en caution, une phrase par partie. Choisissez ses arguments chapitre par chapitre, et en restant dans les limites de l'énoncé.

\end{piege}

\begin{piege}

Sur Verne, vérifiez la fin : le roman ne dit pas que Nemo meurt. Aronnax, rescapé, ignore si le \emph{Nautilus} a résisté au Maelström et souhaite que Nemo vive. Écrire comme un fait que Nemo « sombre avec son secret » force le dénouement. L'exemple de composition du rapport 2026 parle pourtant de la disparition de Nemo dans le maelström : c'est une lecture possible, à condition de la présenter comme une interprétation d'une fin laissée ouverte.

\end{piege}

\hypertarget{deuxiuxe8me-phase-ruxe9diger-le-devoir}{%
\section{Deuxième phase --- Rédiger le devoir}\label{deuxiuxe8me-phase-ruxe9diger-le-devoir}}

Les cinq étapes qui suivent donnent, mises bout à bout, le devoir complet. Sur la copie, aucun titre : seuls les alinéas et un saut de ligne entre les parties marquent la structure.

\hypertarget{uxe9tape-4-ruxe9diger-lintroduction}{%
\subsection{Étape 4 --- Rédiger l'introduction}\label{uxe9tape-4-ruxe9diger-lintroduction}}

\begin{indication}{1}

Cinq moments : une amorce qui situe la citation, la citation recopiée exactement, son analyse rapide, la problématique, l'annonce du plan.

\end{indication}

\begin{indication}{2}

L'amorce peut venir du contexte même de l'article (une forêt menacée en 1872) : c'est plus juste qu'une généralité sur « l'homme et la nature depuis toujours ».

\end{indication}

\hypertarget{solution-3}{%
\subsubsection{Solution}\label{solution-3}}

En 1872, dans un article du \emph{Temps} consacré à la forêt de Fontainebleau, George Sand ne décrit pas un paysage : elle plaide pour une forêt, à une époque où ses futaies sont exploitées et où les peintres ont obtenu que certaines parcelles soient soustraites à la coupe. Le ton est celui d'une prophétesse plus que d'une promeneuse : « Quelle conquête à entreprendre pour l'homme, et je dis pour tout homme actuellement vivant ou à naître ! Entrer dans la nature, chercher l'oracle de la forêt sacrée et rapporter le mot, ne fût-ce qu'un mot, qui doit répandre sur toute sa vie le charme profond de la possession de son être ! » La phrase dessine un trajet en trois gestes, entrer, chercher, rapporter, et le couronne par un état, la possession de soi. Deux lexiques s'y croisent : celui de la conquête, qui est celui de la prise, et celui de l'oracle, qui est celui du sacré. Mais le plus remarquable est le déplacement qu'elle opère en chemin : la conquête annoncée n'a pas pour objet la nature ; ce que l'homme doit finalement posséder, c'est « son être ». Encore ce butin tient-il dans « un mot », et la restrictive « ne fût-ce qu'un mot » tempère d'avance l'enthousiasme de l'exclamation.

Ce programme suppose pourtant beaucoup : qu'il existe un seuil de la nature, que l'homme se tienne d'ordinaire au dehors et puisse y entrer ; que la nature ait quelque chose à dire ; que celui qui l'a entendue puisse revenir et transmettre. Or les trois œuvres au programme racontent ou pensent des immersions. Dans \emph{Vingt mille lieues sous les mers}, Jules Verne fait descendre le professeur Aronnax dans les profondeurs à bord du \emph{Nautilus} du capitaine Nemo ; dans \emph{Le Mur invisible}, Marlen Haushofer enferme sa narratrice, seule avec quelques bêtes, dans une vallée des Alpes autrichiennes que ferme une paroi transparente ; dans \emph{La Connaissance de la vie}, Georges Canguilhem se demande à quelles conditions la pensée peut connaître le vivant sans le manquer. On se demandera donc dans quelle mesure l'expérience de la nature, telle que ces œuvres la donnent à lire, est la conquête dont l'homme reviendrait en possession de lui-même, ou si elle ne vérifie la phrase de Sand qu'à condition d'en renverser le vocabulaire.

Nous verrons d'abord que les trois œuvres suivent bien le trajet que trace Sand : on ne connaît la nature qu'en y entrant, en l'interrogeant et en écrivant ce qu'on en a appris. Nous montrerons ensuite qu'elles récusent les mots par lesquels Sand décrit ce trajet : la conquête y tourne à la prédation, l'oracle y reste muet, le retour y est incertain. Nous établirons enfin que la formule retrouve sa justesse si on la lit à partir de son dernier terme : la possession de son être n'est pas une prise, mais la capacité, pour un vivant parmi d'autres, de se donner des normes et de répondre de ce qui l'entoure.

\begin{methode}

L'analyse de la citation, dans l'introduction, doit déjà produire la tension qui fera avancer le devoir : ici, le décalage entre l'objet annoncé de la conquête et son objet réel, et les présupposés du trajet (un seuil, une parole, un retour).

\end{methode}

\begin{jury}

Rapport 2026 : l'analyse arrive entre la reprise du sujet, « qui, nous le rappelons encore, doit être intégralement reproduit dans l'introduction », et la problématique ; une paraphrase de deux ou trois lignes ne suffit pas. « En un mot : l'analyse du sujet est le cœur de l'introduction. » Le jury relève aussi des fautes de recopiage de la citation et des fautes sur le nom des auteurs.

\end{jury}

\begin{piege}

Les accroches que le rapport 2026 épingle : commencer « brutalement avec le recopiage de la citation », annoncer « La citation que nous allons étudier aujourd'hui », ou partir d'une généralité de développement personnel ou de la conquête spatiale. L'accroche doit déjà montrer une référence en rapport avec le sujet.

\end{piege}

\hypertarget{uxe9tape-5-ruxe9diger-la-premiuxe8re-partie}{%
\subsection{Étape 5 --- Rédiger la première partie}\label{uxe9tape-5-ruxe9diger-la-premiuxe8re-partie}}

\begin{indication}{1}

Lisez la citation de la façon la plus généreuse : chaque sous-partie suit un geste de Sand, et chaque œuvre y apporte une scène ou une thèse précise.

\end{indication}

\begin{indication}{2}

Chez Verne, une scène répond presque mot pour mot à « entrer dans la nature » et à « forêt » : cherchez-la parmi les sorties du \emph{Nautilus}.

\end{indication}

\hypertarget{solution-4}{%
\subsubsection{Solution}\label{solution-4}}

Le premier geste que réclame Sand, entrer, est celui dont les trois œuvres font la condition de toute connaissance. Aronnax est d'abord un savant de cabinet : professeur suppléant au Muséum de Paris, il a publié un ouvrage en deux volumes sur les mystères des grands fonds sous-marins, qu'il n'a jamais vus. Le roman consiste à faire passer ce savoir du livre à la chose, et Verne donne à cette entrée une figure littérale : la promenade en scaphandre dans la forêt sous-marine de l'île Crespo. C'est bien une forêt, et le naturaliste y entre à pied ; il y voit de ses yeux ce que ses livres ne lui avaient pas appris, une végétation qui monte droit vers la surface. Canguilhem formule la même exigence dans l'ordre de la méthode. L'étude consacrée à l'expérimentation en biologie animale montre que le vivant ne se laisse pas traiter comme un corps physique : il est un individu, une totalité, une histoire, et l'on ne peut isoler une fonction, ni conclure d'une espèce à une autre, sans altérer ce qu'on prétend connaître. Le laboratoire lui-même est un milieu, avec ses contraintes propres. Pour connaître le vivant, il faut donc entrer dans les conditions de sa vie au lieu de l'en extraire. Chez Haushofer, enfin, l'entrée n'est pas choisie, elle est subie : venue passer quelques jours au pavillon de chasse de sa cousine Luise et de son mari Hugo, la narratrice découvre au matin que ses hôtes ne sont pas revenus du village et qu'une paroi invisible ferme la route. Le seuil de Sand est devenu un mur qui ne se franchit que dans un sens, mais l'entrée est totale : une citadine doit apprendre la nature par le corps.

Cette entrée ouvre une recherche, et c'est le second geste de Sand : chercher l'oracle, c'est-à-dire interroger une nature qui ne livre pas d'elle-même son sens. Le roman de Verne s'ouvre sur une énigme, ce prétendu monstre marin que les navires croisent en 1866, et il en multiplie d'autres, des ruines de l'Atlantide aux poulpes géants que la science tenait pour fabuleux. Il oppose surtout deux manières de connaître : Conseil, le domestique flamand, range tout poisson dans sa classe et son ordre sans savoir le reconnaître, tandis que Ned Land, le harponneur, le reconnaît sans savoir le classer ; Aronnax note qu'à eux deux ils feraient un naturaliste. Interroger la nature, c'est ajuster la nomenclature à l'expérience, et d'abord chercher le mot juste : quand Ned demande ce qu'est une perle, Aronnax lui donne tour à tour le mot du poète, celui du chimiste et celui du naturaliste, comme si aucun nom ne suffisait seul à dire la chose. Canguilhem donne à cette recherche sa portée philosophique : le vivant pose à la pensée un problème que ne pose pas la matière inerte, parce qu'il est lui-même un centre d'activité qui organise son milieu. Le chercheur ne trouve donc pas seulement des faits, il rencontre des allures, des préférences, des valeurs. La narratrice d'Haushofer, elle, apprend à déchiffrer des signes vitaux : l'état du ciel avant la fenaison, l'humeur de sa vache Bella, le comportement de son chien Lynx, l'approche de la neige. Son oracle est quotidien, mais il est bien une parole à interpréter, et sa survie dépend de la justesse de l'interprétation.

Reste à « rapporter le mot », et les trois œuvres sont, à la lettre, des rapports. Le roman de Verne se donne pour le récit qu'Aronnax rédige après son évasion ; \emph{Le Mur invisible} est le texte que la narratrice écrit sur ce qu'elle appelle elle-même son « rapport », pour ne pas perdre la raison ; \emph{La Connaissance de la vie} est ce qu'un philosophe rapporte de sa fréquentation des biologistes. Et ce mot rapporté revient bien, comme le veut Sand, sur celui qui le rapporte. La narratrice, dont le texte ne donne jamais le nom, se découvre étrangère à la femme qu'elle était à la ville, à ses soucis et à son apparence ; Aronnax, embarqué pour détruire un monstre, finit partagé entre l'envie de fuir et le désir de poursuivre l'exploration. Canguilhem fait de ce retour sur soi la condition même du savoir biologique : la pensée qui connaît la vie est le fait d'un vivant, et, selon la formule de son introduction, « la pensée du vivant doit tenir du vivant l'idée du vivant ». Connaître la vie, c'est aussi se connaître comme vivant.

\begin{methode}

Une référence nommée vaut mieux qu'une allusion : « la promenade en scaphandre dans la forêt de l'île Crespo » démontre une lecture, « une plongée » ne démontre rien.

\end{methode}

\begin{jury}

Rapport 2026 : une sous-partie est un paragraphe, et l'on ne revient pas à la ligne à chaque nouvel auteur, car « les différents auteurs viennent appuyer un même argument dont l'unité se matérialise par le paragraphe que constitue la sous-partie ». Chaque sous-partie reprend les termes du sujet, mais en les analysant : « Le jury n'est pas dupe, et certains candidats reprennent sans cesse les termes du sujet pour dissimuler un manque d'analyse. »

\end{jury}

\hypertarget{uxe9tape-6-ruxe9diger-la-deuxiuxe8me-partie}{%
\subsection{Étape 6 --- Rédiger la deuxième partie}\label{uxe9tape-6-ruxe9diger-la-deuxiuxe8me-partie}}

\begin{indication}{1}

La deuxième partie ne nie pas la première : elle éprouve le vocabulaire de la citation. Prenez un mot par sous-partie (conquête, oracle, rapporter).

\end{indication}

\begin{indication}{2}

Pour chaque mot, montrez le mécanisme de l'échec : conquête punie ou retournée, nature sans sens en soi, retour impossible ou incertain.

\end{indication}

\begin{indication}{3}

Canguilhem fournit l'arrière-plan théorique de chaque objection : la critique du modèle mécanique, la notion de milieu propre, l'écart entre la vie et la connaissance.

\end{indication}

\hypertarget{solution-5}{%
\subsubsection{Solution}\label{solution-5}}

Les trois œuvres suivent donc le trajet de Sand. Mais le suivent-elles dans les termes où elle le décrit ? Aronnax est un prisonnier, la narratrice une captive, l'expérimentateur un homme qui risque d'abîmer ce qu'il veut connaître : c'est le vocabulaire de la phrase, plus que son mouvement, que les œuvres mettent en cause.

Sand parle de « conquête » et de « possession ». Elle ne prône pas la domination de la nature, puisque sa conquête a pour objet « son être » ; mais le mot vient de la guerre, et l'on entend sous lui le programme cartésien qui invitait les hommes à se rendre « comme maîtres et possesseurs de la nature ». Or ce vocabulaire de la prise, dès qu'il s'applique à la nature, les œuvres l'instruisent à charge. Le geste conquérant par excellence est chez Verne aussitôt démenti : lorsque Nemo plante au pôle Sud son pavillon noir et prend possession, en son propre nom, de terres qu'aucun homme n'habite, le \emph{Nautilus} se trouve peu après bloqué sous la banquise par le basculement d'un iceberg, et l'équipage manque mourir asphyxié. Le roman est aussi un livre de tueries : Nemo, qui interdit à Ned Land de harponner des baleines inoffensives pour le seul plaisir et lui reproche les ravages des baleiniers, massacre peu après une troupe de cachalots qu'il juge nuisibles ; et le proscrit qui défendait les opprimés finit par éperonner un navire de guerre et regarder couler son équipage. La conquête rend Nemo semblable à ce qu'il hait. Chez Haushofer, le procès est plus radical encore. Le seul être humain vivant que rencontre la narratrice est un homme qui surgit à l'alpage et tue à coups de hache le jeune taureau et le chien ; elle l'abat d'un coup de fusil. Cet inconnu sans nom ni parole est la prise à l'état pur, détruire sans raison. Le roman ne propose pas pour autant une innocence : la narratrice tue elle aussi, chasse le gibier pour manger et le fait à contrecœur. Canguilhem, enfin, ôte à la maîtrise son fondement. Dans « Machine et organisme », il renverse le modèle cartésien : loin que l'organisme s'explique par la machine, c'est la machine qui s'explique par l'organisme, la technique prolongeant une activité vitale. On ne possède pas comme un objet démontable ce dont on procède.

Sand suppose ensuite une « forêt sacrée » qui rend des oracles. Or la nature, dans les trois œuvres, ne parle pas. \emph{Le Mur invisible} est un roman de montagne qui refuse le lyrisme de la montagne : la beauté d'un matin d'été ne console de rien, des bêtes meurent sans raison, la corneille blanche que la narratrice observe est rejetée par les autres corneilles parce qu'elle est blanche. Et le mur, événement central du livre, n'est jamais expliqué : ni cause, ni auteur, ni sens. Canguilhem donne à ce silence son explication théorique. Dans « Le vivant et son milieu », il reprend à von Uexküll la distinction entre le milieu propre d'un animal, fait des seuls signaux qui ont pour lui une signification, et l'environnement que décrit le physicien : la tique ne vit pas dans le même monde que le savant qui l'observe. Il n'y a donc pas de sens en soi de la nature, qu'un oracle délivrerait ; le sens naît du rapport qu'un vivant institue avec ce qui l'entoure. Chez Verne, enfin, le seul « mot » que le voyage fasse apparaître sous la mer est humain : dans les ruines de l'Atlantide, Nemo trace sur un rocher, avec un morceau de pierre crayeuse, le mot « ATLANTIDE ». Ce n'est pas la mer qui parle, c'est un homme qui écrit sur elle. Et le livre se clôt sur la question de l'Ecclésiaste, demandant qui a jamais pu sonder les profondeurs de l'abîme : l'oracle n'a pas répondu.

Le troisième geste, « rapporter », est peut-être le plus fragile, car il suppose qu'on revienne et qu'on ait quelqu'un à qui parler. La narratrice d'Haushofer ne revient pas : elle écrit au pavillon un rapport que personne ne lira peut-être, et qui s'arrête non parce que son histoire est achevée, mais parce qu'elle n'a plus de papier. Aronnax revient, mais il sait que son récit est invérifiable et se demande lui-même si on le croira. Quant à Nemo, qui détient le savoir le plus complet du voyage, il ne rapporte rien aux hommes : il a confié l'histoire de sa vie à un manuscrit que le dernier survivant de l'équipage doit jeter à la mer, et le roman s'achève sans que l'on sache si le \emph{Nautilus} a résisté au Maelström. Le mot reste au fond. Canguilhem, enfin, met en doute le « charme » même du mot rapporté : son introduction rappelle que connaître, c'est analyser, décomposer ce que la vie compose ; entre le geste de la vie, qui forme, et celui du savoir, qui sépare, il y a un écart que le savoir ne peut abolir. Ce que le savant rapporte n'est jamais la vie elle-même.

\begin{piege}

Faire de Canguilhem un adversaire de la science ou un chantre de la nature contre la technique. Il analyse les conditions de la connaissance du vivant, et « Machine et organisme » fait de la technique un phénomène vital : l'enrôler dans un réquisitoire contre la machine est un contresens.

\end{piege}

\begin{piege}

Transformer la partie en sermon (« la conquête, c'est mal ») et oublier que Sand elle-même écrivait pour défendre une forêt : l'objection porte sur son vocabulaire, non sur son intention.

\end{piege}

\begin{jury}

Rapport 2026 : « Nous ne saurions trop rappeler que l'antithèse n'est pas le contre-pied de la thèse. » Le jury cite une copie qui écrivait d'abord que la conquête de la nature permet la possession de soi, puis, quelques lignes plus loin, qu'elle ne la permet pas. Ici, la deuxième partie ne contredit pas la première : la première vérifie le mouvement de la phrase, la deuxième éprouve son vocabulaire.

\end{jury}

\hypertarget{uxe9tape-7-ruxe9diger-la-troisiuxe8me-partie}{%
\subsection{Étape 7 --- Rédiger la troisième partie}\label{uxe9tape-7-ruxe9diger-la-troisiuxe8me-partie}}

\begin{indication}{1}

Ne conciliez pas mollement : relisez la citation jusqu'à son dernier mot. Qu'est-ce qui change si l'objet de la possession est soi-même ?

\end{indication}

\begin{indication}{2}

Cherchez chez Canguilhem un concept qui donne un sens non propriétaire à « posséder son être ».

\end{indication}

\begin{indication}{3}

Chez Verne, comparez les deux hommes qui ont fait le même voyage : lequel possède la mer, lequel a changé ?

\end{indication}

\hypertarget{solution-6}{%
\subsubsection{Solution}\label{solution-6}}

Faut-il alors conclure que la phrase de Sand n'éclaire rien, et qu'elle appartient à un siècle que nos œuvres auraient corrigé ? Ce serait manquer ce qu'elle a de juste et qui n'apparaît qu'à la relire jusqu'à son dernier mot : ce que Sand promet n'est pas la possession de la nature, mais celle de « son être ».

Il faut d'abord renoncer à l'image du seuil. On n'entre pas dans la nature, parce qu'on n'en est jamais sorti. C'est la leçon la plus forte de Canguilhem : l'homme de science peut construire un univers sans centre ni lieu privilégié, il n'en reste pas moins un vivant pour qui le milieu s'organise autour de lui ; la pensée qui prétend connaître la vie est une production de la vie. Il n'y a donc pas de porte à pousser, mais une séparation, instituée par le savoir lui-même, qu'il faut reconnaître. Haushofer donne à cette idée sa figure la plus saisissante : en enfermant la narratrice, le mur lui révèle qu'elle était déjà dedans. Soumise au froid, à la faim et à la fatigue comme sa vache et son chien, elle se découvre une créature parmi d'autres ; et sur l'alpage, dans le silence de la prairie, elle éprouve qu'il lui devient presque impossible de rester un moi séparé du reste du vivant. Verne, de son côté, montre que l'immersion passe par des machines : le scaphandre, les réservoirs d'air, le \emph{Nautilus} lui-même, dont la devise, \emph{Mobilis in mobili}, dit un être mobile dans l'élément mobile. La machine n'y est pas l'ennemie de la nature ; elle est, au sens de Canguilhem, le prolongement d'un vivant qui étend son milieu. Les grandes vitres du salon ne séparent pas seulement Aronnax de l'océan : elles le lui donnent à voir.

Il faut ensuite prendre au sérieux la modestie de Sand : elle n'annonce pas un secret, elle annonce « un mot », et encore « ne fût-ce qu'un mot ». Or c'est précisément ce que rapportent nos trois œuvres, qui sont elles-mêmes ce mot. Aronnax écrit un récit qu'il sait incomplet, mais il peut affirmer qu'à la question de l'Ecclésiaste deux hommes seulement ont désormais le droit de répondre, Nemo et lui. La narratrice d'Haushofer écrit sans style, sans lecteur assuré, presque sans espoir ; mais tant qu'elle écrit, elle demeure quelqu'un, et son rapport lui permet de mettre de l'ordre dans ce qu'elle a vécu, jusqu'à la mort de Lynx et du taureau, vers laquelle tout son récit s'achemine. Canguilhem, lui, ne rapporte pas de la biologie des résultats, mais une manière de penser : une philosophie qui accepte d'apprendre du vivant au lieu de lui imposer ses catégories. Le mot rapporté ne délivre pas le sens de la nature ; il maintient vivant celui qui l'écrit et le rend transmissible. C'est un usage plus modeste que l'oracle de Fontainebleau, et plus solide.

Il faut enfin donner à « la possession de son être » son sens véritable. Posséder son être, ce n'est pas s'appartenir comme on possède un bien ; Canguilhem fournit le concept qui convient. La vie, selon lui, est une activité normative : être en bonne santé, ce n'est pas se conformer à une moyenne, c'est pouvoir instituer de nouvelles normes quand le milieu change, pouvoir tomber malade et s'en relever. Se posséder, en ce sens, ce n'est pas se tenir, c'est pouvoir changer. Telle est exactement la conquête de la narratrice d'Haushofer. Les normes sociales ont disparu avec les hommes ; elle en institue d'autres, un calendrier des travaux, une répartition des réserves, une règle de conduite envers les bêtes, et ce sont ces normes qui font d'elle un sujet plutôt qu'une simple survivante. Il n'y a chez elle aucun « charme profond », aucune extase dans la forêt sacrée ; il y a le soin, qui est devenu sa manière d'être. Verne fournit le contre-exemple et, par là, la confirmation. Nemo possède la mer, l'indépendance, un empire où il ne reconnaît aucun maître ; mais il ne se possède pas, car sa haine décide de sa route jusqu'à l'hécatombe, et la nuit de la fuite Aronnax l'entend implorer Dieu de faire cesser ce qu'il a lui-même déchaîné. Aronnax, lui, ne possède rien et rapporte un récit lacunaire : c'est pourtant lui qui a changé. Le roman sépare ainsi en deux personnages ce que la phrase de Sand risquait de confondre, la conquête de la nature et la possession de soi.

\begin{methode}

Une bonne troisième partie retourne les objections de la deuxième en conditions de vérité de la citation : la modestie du « mot » et le déplacement vers « son être » étaient déjà dans la phrase.

\end{methode}

\begin{jury}

Rapport 2026 : « La troisième partie est souvent, même quand les copies sont de bonne tenue, la moins réussie. » Le jury attend un dépassement, ni mélange des deux premières parties ni retour à la première : « Dépassement, oui, mais jamais effacement du sujet derrière une question beaucoup trop vague qui aurait pu servir de troisième partie à n'importe quel sujet. »

\end{jury}

\begin{piege}

Les troisièmes parties qui, selon le jury, « sentaient le réchauffé » : l'adaptation de l'homme, la nature qui reprend toujours ses droits, l'homme produit de la nature. La sous-partie sur le seuil, ci-dessus, frôle ce dernier thème ; il ne vaut que parce qu'il est greffé sur un mot du sujet (« entrer ») et qu'il en tire une conséquence (il n'y a pas de seuil à franchir).

\end{piege}

\hypertarget{uxe9tape-8-ruxe9diger-la-conclusion}{%
\subsection{Étape 8 --- Rédiger la conclusion}\label{uxe9tape-8-ruxe9diger-la-conclusion}}

\begin{indication}{1}

Répondez à la question posée en mesurant : sur quoi la citation éclaire, sur quoi elle cesse d'éclairer.

\end{indication}

\begin{indication}{2}

Pour l'ouverture, regardez si un segment de la citation n'a pas encore été exploité.

\end{indication}

\hypertarget{solution-7}{%
\subsubsection{Solution}\label{solution-7}}

La phrase de George Sand éclaire donc nos trois œuvres par son point de départ et par son terme plus que par le chemin qu'elle décrit entre les deux. Elle a raison de poser qu'il faut entrer dans la nature pour la connaître : Verne, Canguilhem et Haushofer refusent tous trois le savoir de cabinet, le vivant isolé de son milieu, le paysage regardé de loin. Elle a raison encore de tenir que ce savoir revient sur celui qui l'acquiert. Mais les mots qu'elle emploie, conquête, oracle, rapporter, sont démentis par les œuvres : la conquête y tourne à la prédation, l'oracle s'y tait parce que la nature n'a pas de sens en soi, et le retour, chez Nemo comme chez la narratrice, reste incertain ou impossible.

Ce qui subsiste n'est pas rien : c'est même l'essentiel. On ne connaît pas le vivant sans en être ; on ne rapporte de la nature qu'un mot, et ce mot, chez nos trois auteurs, est un livre ; surtout, la conquête annoncée n'était pas celle de la nature. Ce que Sand nomme la possession de son être, nos œuvres l'appellent autrement : la normativité du vivant chez Canguilhem, le soin et la règle que s'invente la narratrice d'Haushofer, la liberté intérieure qui manque à Nemo, maître de la mer et esclave de sa vengeance.

Reste la clause que Sand jette dans une incise et qui, en 1872, engageait le sort d'une forêt réelle : « pour tout homme actuellement vivant ou à naître ». Si l'expérience de la nature est due aux hommes à naître, alors elle oblige ceux qui vivent, et la conquête à entreprendre se renverse en une responsabilité à assumer.

\begin{piege}

L'ouverture convenue sur l'urgence écologique plaque sur les œuvres un discours qu'elles ne tiennent pas. Prise dans la citation elle-même (« vivant ou à naître »), l'ouverture prolonge le raisonnement au lieu de le décorer.

\end{piege}

\begin{jury}

Rapport 2026 : « Le jury rappelle que la conclusion n'est pas un résumé, mais un bilan. » Il reproche les conclusions qui redisent le plan au passé composé (« Nous avons vu que\ldots{} ») au lieu de faire saillir, avec les termes de la problématique, les lignes de force de la confrontation entre les œuvres et la citation.

\end{jury}

\hypertarget{bilan}{%
\section{Bilan}\label{bilan}}

\begin{itemize}
\tightlist
\item
  \textbf{Méthode de la dissertation sur citation.} Analyser la forme avant le sens ; séparer ce que la phrase affirme de ce qu'elle présuppose ; problématiser sur une tension interne (ici, mouvement juste, vocabulaire démenti) ; indexer chaque partie sur des termes de la citation.
\item
  \textbf{Confrontation des œuvres.} Chaque sous-partie fait dialoguer au moins deux œuvres, et chaque œuvre sert dans chaque partie. Le motif du « rapport » (le récit d'Aronnax, le « rapport » de la narratrice, le livre du philosophe) est un fil transversal précieux pour tout sujet sur la connaissance ou la transmission de l'expérience.
\item
  \textbf{Canguilhem à maîtriser.} La pensée du vivant procède du vivant ; le vivant est une totalité et l'expérimentation doit en tenir compte ; la machine s'explique par l'organisme ; le milieu propre (von Uexküll) ; la santé comme normativité. Ces cinq idées s'adaptent à presque tous les sujets du thème.
\item
  \textbf{Ce qui a départagé les copies en 2026.} Selon le rapport, l'analyse de la citation entière (et non du seul mot « conquête »), une problématique unique et une troisième partie qui dépasse sans quitter le sujet. Le jury conclut que trop de candidats étaient « des lecteurs de fiches, de corrigés tout rédigés, de recueils de citations, et non des lecteurs d'œuvres » : un corrigé comme celui-ci se lit pour la méthode, non pour en réciter les exemples.
\item
  \textbf{Exactitude.} Les noms (Aronnax, Conseil, Ned Land, Nemo ; Lynx et Bella), la devise \emph{Mobilis in mobili}, le dénouement ouvert de Verne, le rapport interrompu faute de papier chez Haushofer : une erreur sur ces points se voit et coûte cher.
\item
  \textbf{Sujets transposables.} Toute citation qui oppose maîtrise et appartenance, savoir et expérience, nature muette et nature signifiante, ou qui interroge la possibilité de transmettre une expérience, se traite avec la même grille : accord sur le mouvement, objection au vocabulaire, relecture par le terme le plus riche.
\end{itemize}

\end{document}
